\section{A direct cycle deficit in a growing girth window}
\label{sec:random-girth-window}

Moore-ball packing makes the shortest-cycle prefix argument quantitative.
The estimate below does not require a restriction on cyclic cuts.
All logarithms are binary.

\begin{lemma}[A direct bound for the large prefix branches]
\label{lem:random-girth-prefix-path}
There are absolute constants $C,R$ with the following property.
Let $r\ge R$, put $x=\log r$, $z=\log x$, and set
\[
 g_0=\left\lceil4z+2\log z+12\right\rceil.
\]
Let $N$ be a normalized $k$-pole with rank $b\in[r/2,r]$, girth at least $g_0$, and $k\le3x$.
For any distinct ports $u,v$,
\[
 P_{uv}(N)\le
 \left(\frac12+C\frac{x^{11/2}}{\sqrt r}\right)2^b.
\]
\end{lemma}

\begin{proof}
The order of $N$ is $n=2b+k-2$.
Apply Lemma~\ref{lem:four-pole-moore-geodesic} with the integer girth lower bound $g_0$.
The terminal distance $\ell$ satisfies
\[
 \ell+1<g_0\left(k+1+n2^{-g_0/2}\right).
\]
For sufficiently large $r$,
\[
 g_0\le8z,\qquad
 2^{-g_0/2}\le\frac1{64x^2z},\qquad
 n\le3b.
\]
Hence
\[
 \ell+2\le8z(3x+1)+\frac{3b}{8x^2}+1
 \le\frac{b}{2x^2}.
\]
The last comparison is uniform for $b\ge r/2$.

Suppress every degree-two vertex other than $u,v$.
This preserves rank and terminal-path count, and it can only shorten terminal distance.
Call the resulting two-pole $N'$.
Parallel edges or loops do not affect these assertions or the multigraph forest theorem.
Add a new two-edge $u$--$v$ path to $N'$, obtaining a subcubic graph $H$ of rank $\rho=b+1$.
Let $A$ be a shortest terminal path in $N'$ and let $W$ be its union with the new path.
Then $W$ is a cycle and
\[
 |V(W)|\le\ell+2\le\frac{b}{2x^2}
 \le\frac{\rho}{(\log\rho)^2}.
\]
For the final inequality, use $\rho\le r+1$ and take $r$ sufficiently large.
Every vertex outside $W$ is cubic.

Each terminal path $Q$ of $N'$ gives two even subgraphs of $H$:
its union with the new path, and $Q\mathbin{\triangle}A$.
These $2P_{uv}(N')$ subgraphs are distinct.
Outside $V(W)$, each is a subgraph of $Q$ and is therefore a forest.
Consequently
\[
 \frac{P_{uv}(N)}{2^b}
 \le
 \Pr\bigl(\xi_H-V(W)\text{ is a forest}\bigr).
\]
The endpoint forest theorem, Theorem~\ref{thm:parallel-links-log-five-half-forest}, bounds the right side by
\[
 \frac12+O\!\left(\frac{(\log\rho)^{11/2}}{\sqrt\rho}\right)
 \le\frac12+C\frac{x^{11/2}}{\sqrt r}.
\]
All constants are absolute.
\end{proof}

\begin{theorem}[The quantitative girth window]
\label{thm:random-girth-window}
There are absolute constants $C,R$ such that every connected loopless cubic multigraph $G$ of rank $r\ge R$ and girth $g$ satisfying
\[
 \left\lceil4\log\log r+2\log\log\log r+12\right\rceil
 \le g
\]
obeys
\[
 c(G)\le
 2^{r-1}-2^{r-g-1}
 +C\,2^r\frac{(\log r)^{11/2}}{\sqrt r}.
\]
No hypothesis on cyclic connectivity is required.
\end{theorem}

\begin{proof}
Choose a shortest cycle $C=v_0\cdots v_{g-1}v_0$.
The cubic Moore bound gives $g\le2\log(2r-2)+2\le3\log r$ for sufficiently large $r$.
Delete its edge $e=v_0v_{g-1}$ and write $b=r-1$.
Partition the remaining terminal paths by first divergence from the fixed path $C-e$.
Lemma~\ref{lem:path-shortest-cycle-prefix} gives, for each $1\le p\le g-1$, either no branch or a normalized pole of rank
\[
 b_p\le b-p.
\]
Its two terminals are distinct ports.
It has at most $g$ ports and girth at least $g$.
The last branch includes the pendant-endpoint deletion stated in that lemma.

Put $\varepsilon_r=C_1(\log r)^{11/2}/\sqrt r$, with $C_1$ from Lemma~\ref{lem:random-girth-prefix-path}.
If $b_p\ge r/2$, that lemma bounds the branch count by
\[
 (1/2+\varepsilon_r)2^{b_p}
 \le(1/2+\varepsilon_r)2^{b-p}.
\]
If $b_p<r/2$, the parity-class bound gives at most $2^{r/2}$ paths.
Including the fixed path $C-e$ and summing the geometric series therefore gives
\begin{align*}
 P_{v_0v_{g-1}}(G-e)
 &\le1+(1/2+\varepsilon_r)\sum_{p=1}^{g-1}2^{b-p}
       +(g-1)2^{r/2}\\
 &=1+(1/2+\varepsilon_r)(2^b-2^{b-g+1})
       +(g-1)2^{r/2}.
\end{align*}

The unrestricted endpoint cycle bound applies to the connected graph $G-e$:
\[
 c(G-e)\le2^{b-1}
 +O\!\left(2^b\frac{(\log b)^{11/2}}{\sqrt b}\right).
\]
Add this estimate to the terminal-path count.
The two half terms give $2^{r-1}-2^{r-g-1}$.
Since $g\le3\log r$, the term $1+(g-1)2^{r/2}$ is absorbed by the displayed endpoint error.
Increasing the absolute constant proves the theorem.
\end{proof}

\begin{corollary}[A strict deficit up to half the rank logarithm]
\label{cor:random-girth-window-strict}
There is an absolute constant $C_0$ such that, for sufficiently large $r$, every connected loopless cubic multigraph of rank $r$ and girth in the interval
\[
 \left\lceil4\log\log r+2\log\log\log r+12\right\rceil
 \le g
 \le\frac12\log r-\frac{11}{2}\log\log r-C_0
\]
satisfies
\[
 c(G)\le2^{r-1}-2^{r-g-2}<2^{r-1}.
\]
More generally, within the theorem's girth range, if
\[
 2^g\frac{(\log r)^{11/2}}{\sqrt r}\longrightarrow0,
\]
then the negative term is $(1-o(1))2^{r-g-1}$.
\end{corollary}

\begin{proof}
Choose $C_0$ so that $C2^{-C_0}\le1/4$, where $C$ is the theorem's error constant.
The upper girth bound then gives
\[
 C\frac{(\log r)^{11/2}}{\sqrt r}
 \le2^{-g-2}.
\]
Subtract this error from $2^{-g-1}$ after normalizing by $2^r$.
The limiting assertion follows from the same comparison.
\end{proof}

The remaining unrestricted problem includes girths between the available double-logarithmic ranges and girths above this window.
This theorem supplies a quantitative exclusion range; it does not prove the unrestricted sharp-order error.
